\section{Auditing a claimed route from Jones identity to recognition}
\label{sec:jones-detection-audit}

A June 2026 preprint by Carmi and Cohen claims that the Jones polynomial
detects the unknot \cite{carmi-cohen2026}.  We checked a specific prerequisite:
Lemma~5.8 bounds residual Jones--Vassiliev degree by
$\widehat B+2$, under its machine hypotheses.  Section~5.1.1 sets
$\widehat B=3n+\ell-1+11$ from the pre-planting host, and
Proposition~5.12 claims the required local construction at a seed crossing,
recovering an ordinary odd anti-parallel twist.  Together these statements
would bound every such residual twist on a three-crossing knot host by
$3\cdot3+1-1+11+2=22$, independently of its length.

The explicit computation below contradicts this combined consequence.
It does not distinguish an error in the conditional algebraic lemma from
an error in the claimed geometric realization, and it does not provide
a counterexample to Jones unknot detection itself.  It is sufficient to
reject this proof as a dependency of the maintained recognition algorithm.
The stronger Jones identity backend therefore continues to require an
independent unknot certificate.

\subsection{An explicit host and a twist with fixed exterior}

Use the oriented trefoil with PD rows
\[
 H=((5,1,2,0),(1,4,3,2),(0,3,4,5)).
\]
Keep its first two rows fixed.  Replace the last row by a linear twist
of odd length $L$, preserving its four exterior labels $0,3,4,5$.
Start with $(a,b)=(0,3)$.  For $i=0,\ldots,L-2$, insert
$(a,b,6+2i,7+2i)$ and replace $(a,b)$ by $(7+2i,6+2i)$.
Finish with $(a,b,4,5)$.  Call the resulting knot diagram $D_L$.
For $L=1$ this is precisely $H$, and $D_L$ has $L+2$ crossings.
All of its crossings are positive in the orientation given by traversing
opposite ports and equal-label arc endpoints.

The two strands in the inserted band are anti-parallel: at every inserted
row, exactly one of ports $0,1$ is incoming and exactly one of ports $2,3$
is incoming.  The oriented smoothing pairs ports $(0,1)$ and $(2,3)$.
After smoothing any one band crossing, the others successively become
adjacent-port monogons.  Deleting these by Reidemeister~I moves gives
exactly the two exterior rows obtained by smoothing the seed of $H$.
All moves take place in the inserted band.  This shadow calculation is
independent of the signs assigned to the remaining crossings.

There is also an explicit host anchor in every odd sign cube: assign the
first $L-1$ crossings alternating signs $+,-$, and the last sign $+$.
Cancel each consecutive opposite-sign pair by Reidemeister~II.  The
surviving row and its exterior are exactly $H$, up to arc relabeling.
Thus the audit does not infer applicability merely from a knot-family
name or from a numerical resemblance of polynomials; it checks the
specified local anti-parallel and anchor behavior.

\subsection{An independent polynomial calculation}

Use $t$ for the standard Jones variable and $z=t^{1/2}$ for the variable
called $x$ in the audited preprint.  This $z$ is different from the
$x=A^4=t^{-1}$ used in our numerical Potts ring.
The seed polynomial is
\[
 V_1(t)=t+t^3-t^4.
\]
Oriented smoothing of a band crossing gives the positive Hopf link,
with $V_0(z)=-z-z^5$.  The four bracket states of the two-crossing
smoothed host independently check this value.  Switching one band
crossing allows a Reidemeister~II cancellation and produces $D_{L-2}$.
The usual oriented Jones skein relation consequently gives
\[
 t^{-1}V_L(t)-t V_{L-2}(t)
 =(t^{1/2}-t^{-1/2})(-t^{1/2}-t^{5/2}),
\]
or
\[
 V_L(t)=t^2V_{L-2}(t)+t-t^2+t^3-t^4\qquad(L\ge3\text{ odd}).
\]
In particular, the leading term is $-t^{L+3}$.  For the smallest positive
example exceeding the claimed cap, $L=9$,
\[
 V_9(t)=t-t^2+2\sum_{j=3}^{9}(-1)^{j+1}t^j-t^{10}+t^{11}-t^{12}.
\]
An independent whole-cube bracket computation checks this eleven-crossing
diagram using only $2^{11}$ states.  It uses neither a Tait graph nor the
production quadratic evaluator.  The exact faithful backend independently
recovers the same complete polynomial.

Now write $p=z-z^{-1}$, so $z^2=pz+1$.  Reducing $V_L(z^2)$ into
$a(p)+b(p)z$ is a formal operation in a free rank-two
$\mathbb Z[p]$-module.  Multiplication by $z$ sends the pair $(a,b)$
to $(b,a+pb)$.  Thus for positive exponent $k$, the coefficient of $z$
in $z^k$ has degree $k-1$ and leading coefficient one.  The term
$-z^{2L+6}$ therefore contributes a unique leading term
$-p^{2L+5}$ to $b(p)$; every other Jones monomial has exponent at most
$2L+4$ and cannot cancel it.  Consequently
\[
 \deg_p V_L=2L+5,\qquad
 f_{23}(D_9)=(a_{23},b_{23})=(0,-1)\ne(0,0).
\]
This contradicts the degree-22 consequence above.  There is no numerical
approximation or degree estimation involved.  Reflection provides another
check: the mirrored family has degree $2L+6$, so the mirror of $D_9$ has
degree 24.  For any fixed proposed cap, the same recurrence supplies
larger counterexamples to a uniform residual-twist bound.

\subsection{Reproducibility and the scope of the finding}

The maintained research audit is
\path{fast/jones_research/audit_residual_twists.py}, with results in
\path{fast/results/residual_twist_audit_20261008.json}.

It checks lengths $1,3,5,7,9,23,31$, both mirrors, and stores each PD,
full Jones polynomial, both formal JVP coordinates and their actual degrees.
Ten small cases have whole-cube comparisons; all fourteen have faithful
recovery and skein-recurrence comparisons.  Formal substitution
$p=z-z^{-1}$ reconstructs every original Laurent polynomial, checking the
coordinate conversion independently of the numerical quadratic ring.
The longer cases include lengths beyond 22, so the discrepancy is not
resolved by requiring a twist longer than that threshold.

The geometry record contains 79 individual smoothings and 1576 explicit
Reidemeister~I deletions with the fixed exterior retained, together with
the alternating-sign Reidemeister~II anchor traces.  Another 42 crossing-sign
assignments, exhausting sign cubes of lengths one, three and five, compare
whole-cube polynomials with their signed-sum twist reductions.  For each
assignment, the independent bracket computation sums all smoothing states.

The audit does not construct the four-clasp machine or identify the precise
step of its claimed collapse that fails.  Its conclusion is the narrower,
fully reproducible one: the uniform residual cap and universal geometric
realization asserted by the audited statements cannot both be correct.
In particular, the preprint does not close our gap between fast Jones
computation and fast complete recognition.
